\chapter{Institutional setting}\label{ch:inst}
\chapterlead{The Indian rules that shape each limit to arbitrage: who may sell short, which stocks have derivatives, how
the position limit and the ban work and how they were enforced, what trading costs, and where arbitrage capital comes
from. Figure~\ref{fig:timeline} places the dated changes on one time line.}

\begin{figure}[t]
\caption[The dated changes to the limits, 2015 to 2026]{Almost every limit to arbitrage changed on a known date during the sample}\label{fig:timeline}
\exsub{Futures STT, enforcement of the ban, the position-limit rule and arbitrage capital, 2015 to 2026}
\centering\includegraphics[width=\linewidth]{fig_timeline}
\exsource{Source: SEBI circulars and press releases, NSE circulars (listed in the references).}
\end{figure}

\section{Short selling and securities lending}
Under SEBI's short-selling framework (restated in a circular of 5 January 2024, in line with the 2007 framework), all
classes of investors may sell short, but naked short selling is prohibited: every short sale must be delivered at
settlement. Institutional investors may not day trade. Securities traded in the F\&O segment are eligible for short
selling. An overnight short in the cash market therefore needs borrowed stock, which comes from the securities lending
and borrowing (SLB) segment. SLB covers F\&O stocks and, since July 2013, other liquid stocks that meet the exchanges'
criteria, but trading in it is thin: on a typical day a few dozen stocks are lent. In practice the overnight short in an
Indian stock is a short single-stock future (or an options position), and this exists only for F\&O stocks. Stock
derivatives have been physically settled since the October 2019 expiry.

\section{The F\&O segment}
A stock enters the F\&O segment when it meets SEBI's eligibility criteria (market capitalisation, traded value, order size,
MWPL and delivery value; the criteria were tightened in 2018 and again on 30 August 2024) and leaves it after failing them.
From 2015 to October 2026 there were 239 entries and 171 exits, 162 of them eligibility exits after which the stock kept
trading in the cash market. Entry gives a stock futures and options and formal short-selling eligibility at once, and
exit removes them.

\section{The market-wide position limit and the ban}
\begin{table}[t]
\caption[Position-limit rules over the sample]{The limit was cut and open interest redefined in October 2025}\label{tab:mwpl}
\exsub{Market-wide position limit (MWPL) and the open interest tested against it}
\small\begin{tabularx}{\textwidth}{@{}K{3.2cm}LL@{}}
\toprule
\thd{Period} & \thd{MWPL} & \thd{Open interest tested} \\
\midrule
To 30 Sep 2025 & 20\% of shares held by non-promoters (free float) & Gross open interest, futures and options \\ \rowrule
20 Mar to 26 Nov 2020 & Halved for volatile or heavily used stocks (61 stocks) & Gross open interest \\ \rowrule
From 1 Oct 2025 & Lower of 15\% of free float and 65 times average daily delivery (in shares), floor 10\% of free float; recomputed quarterly & Futures-equivalent (delta-adjusted) open interest \\
\bottomrule
\end{tabularx}
\exsource{Source: SEBI circulars of 22 September 2006 and 29 May 2025; SEBI press release 18/2020; NSE Clearing circular 119/2025.}
\end{table}

A stock enters the ban when market-wide open interest exceeds 95\% of its MWPL at the end of a day, from the next day, and
leaves it when open interest falls to 80\%. During the ban every participant may only reduce positions, in all expiries
and in both futures and options; rolling a position into the next expiry is not permitted. How the rule was enforced
changed twice, and the changes matter for the results:
\begin{itemize}
\item \emph{Until 22 March 2020}, violations were judged on end-of-day positions, so positions opened and closed within a
day were not breaches.
\item \emph{From 23 March 2020}, after SEBI's press release 18/2020, the clearing corporations judged violations on peak
intraday positions; this was kept ``till further directions'' when the other March 2020 measures were withdrawn in November 2020.
\item \emph{From 1 October 2025}, under SEBI's circular of 29 May 2025 and NSE Clearing circular 119/2025, an entity's
end-of-day futures-equivalent position must fall relative to its position when the ban began. This is more permissive
than the earlier rule of trading only through offsetting positions: new contracts that lower the delta (for example a long
put against a long future) are allowed. The package as a whole was not a relaxation: every stock's MWPL fell, and
intraday monitoring of MWPL utilisation began on 3 November 2025.
\end{itemize}
Increases in a banned stock are penalised (from October 2025: 1\% of the value of the increase, minimum Rs 5,000 and
maximum Rs 1 lakh per entity per stock per day).

\section{Costs}
The STT on the sale of futures was 0.01\% from June 2013, 0.0125\% from 1 April 2023, 0.02\% from 1 October 2024 and
0.05\% from 1 April 2026. Delivery trades in the cash market pay 0.1\% on each side, and physical settlement of stock
derivatives pays 0.1\%. Stamp duty on purchases has been uniform since July 2020 (0.015\% for delivery, 0.002\% for futures).
Exchange transaction charges were revised on 1 October 2024.

\section{Arbitrage capital}
Arbitrage mutual funds hold long cash, short futures positions and are taxed as equity funds. When the taxation of debt
funds changed on 1 April 2023 (the same day as an STT increase), arbitrage funds became the tax-efficient home for
short-term money. Their net assets were Rs 56,218 crore in April 2018, Rs 71,106 crore in April 2023, Rs 1.99 lakh crore in
October 2024 and Rs 2.98 lakh crore in September 2026 (Figure~\ref{fig:aum}).

\begin{figure}[t]
\caption[Net assets of arbitrage mutual funds, 2018 to 2026]{Arbitrage-fund assets quadrupled after debt funds lost their tax advantage}\label{fig:aum}
\exsub{Monthly net assets of the arbitrage-fund category, April 2018 to September 2026}
\centering\includegraphics[width=\linewidth]{fig_arb_aum}
\exsource{Source: AMFI monthly reports; author's calculations.}
\end{figure}

\section{Prices}
NSE sets a stock future's daily settlement price at a theoretical cost-of-carry value when the contract does not trade in
the last half hour. A settlement-price basis would therefore shrink mechanically on thin days, including ban days; all
measures in this thesis use traded prices. From 3 August 2026 the closing price of F\&O stocks comes from a closing
auction.
